\chapter{Installation}
\label{ch:install}

\section{Requirements}

\begin{itemize}
  \item Python 3.9 or newer;
  \item NumPy 1.22 or newer and SciPy 1.8 or newer;
  \item \pyscf{} 2.10 or newer;
  \item optional: \code{mpi4py} 3.1 or newer, for the parallel $k$-point code
        (an MPI library such as OpenMPI must be installed first);
  \item optional, for development: \code{pytest}, \code{ruff} and
        \code{pre-commit}.
\end{itemize}

The dependencies are declared in \code{pyproject.toml}. \code{pip} installs
NumPy, SciPy and \pyscf{} automatically if they are missing.

\section{Installing from the repository}

It is recommended to install \pycmf{} into its own virtual environment:

\begin{lstlisting}[style=shell]
$ git clone https://github.com/QuantumLab-HCMUS/pyCMF-v1.git
$ cd pyCMF-v1
$ python3 -m venv .venv
$ source .venv/bin/activate
$ pip install --upgrade pip
$ pip install .                # users
$ pip install -e ".[test]"     # editable install with pytest
$ pip install -e ".[dev]"      # also ruff and pre-commit
$ pip install -e ".[mpi]"      # also mpi4py
\end{lstlisting}

With an editable install (\code{-e}) changes to the source code in
\code{src/pycmf} take effect without reinstalling.

\subsection*{Another PySCF on PYTHONPATH}

If the environment variable \code{PYTHONPATH} points to another checkout of
\pyscf{} (for example an older fork), that checkout takes precedence over the
\pyscf{} installed in the virtual environment. \pycmf{} then fails or, worse,
runs with the wrong version. Unset the variable before working with
\pycmf{}:

\begin{lstlisting}[style=shell]
$ unset PYTHONPATH
$ source .venv/bin/activate
$ python -c "import pyscf; print(pyscf.__version__, pyscf.__file__)"
\end{lstlisting}

The last command must print a version of at least 2.10 and a path inside
\code{.venv}. A single command can also be run with the variable removed,
\code{env -u PYTHONPATH .venv/bin/python script.py}.

\section{Checking the installation}

\begin{lstlisting}[style=shell]
$ python -c "import pycmf; print(pycmf.__version__)"
1.0.0
$ pytest -m "not slow"      # fast tests
$ pytest                    # all tests, including the slow k-point tests
\end{lstlisting}

The fast part of the suite has 106 tests and runs in a few minutes; the full
suite adds nine slower $k$-point tests (symmetry, larger meshes, MPI). Without
\code{mpi4py} the MPI test is skipped. A complete run on a 20-core
workstation took about 7~minutes for the fast tests and 2~minutes for the slow
ones.

The examples can be compared with their reference output:

\begin{lstlisting}[style=shell]
$ cd examples
$ python check_examples.py          # stable examples
$ python check_examples.py --all    # also the experimental ones
\end{lstlisting}

\section{Parallel execution}

\pycmf{} uses two levels of parallelism.

\begin{description}
  \item[Threads.] All large tensor contractions are matrix products that call
        the BLAS library of NumPy and the multithreaded routines of \pyscf{}.
        The number of threads is set with \code{OMP_NUM_THREADS}.
  \item[MPI.] The periodic code (\code{KOBMP2}) distributes the loop over
        $k$-points with \code{mpi4py}. The same script runs on every
        process; see Section~\ref{sec:kpt-mpi}.
\end{description}

The product of MPI processes and threads per process must not exceed the
number of physical cores. Oversubscribed BLAS threads slow a calculation down
by orders of magnitude. MPI can be switched off without uninstalling
\code{mpi4py} by setting the environment variable \code{PYCMF_NO_MPI=1}.

\section{Memory and scratch space}

The memory limit of a calculation is the attribute \code{max_memory} (in MB),
which is taken from the mean-field object. Large intermediate arrays are
written to the scratch directory of \pyscf{} (\code{TMPDIR} or
\code{pyscf.lib.param.TMPDIR}). Periodic calculations with density fitting
can need several gigabytes of disk space for the three-index integrals; see
Section~\ref{sec:kpt-resources}.

\section{Quick start}

The following script computes the OBMP2 energy of the OH radical and shows
the main stable methods.

\begin{lstlisting}[style=python]
from pyscf import gto, scf
import pycmf

mol = gto.M(atom='O 0 0 0; H 0 0 0.97', basis='cc-pvdz', spin=1)
mf = scf.UHF(mol).run()

mp = pycmf.OBMP2(mf).run()            # restricted or unrestricted, from mf
print(mp.e_tot, mp.spin_square())

o2 = pycmf.O2BMP2(mf, c_os=1.2).run() # c_os has no default

mf_df = scf.UHF(mol).density_fit().run()
dh = pycmf.OBDH(mf_df).run()          # B2PLYP by default

ham = pycmf.OBDF(mp, ncas=5, nelecas=(3, 2))   # active-space Hamiltonian
e_tot = ham.kernel()                           # FCI in the active space
h1, e_core = ham.get_h1eff()
h2 = ham.get_h2eff()
\end{lstlisting}
